\documentclass[11pt]{article}

\usepackage[margin=1in]{geometry}
\usepackage{booktabs}
\usepackage{amsmath}
\usepackage{microtype}
\usepackage{xcolor}
\usepackage[hidelinks]{hyperref}
\usepackage{caption}
\captionsetup{font=small, labelfont=bf}

\title{Schemas in Weights, Facts in the Organizer:\\
A Controlled Test of Fact Offloading in Small Language Model Pretraining}

\author{Stephen Zhang\\
\small Memory Split Project, July 2026\\
\small Preliminary results paper for internal review; confirmatory battery in progress}

\date{July 20, 2026}

\begin{document}
\maketitle

\begin{abstract}
Small language models are widely believed to waste parametric capacity on
factual memorization. The phi-3 team filters facts out of training data
explicitly to preserve capacity for reasoning, and recent memory-augmented
architectures are motivated by the same intuition. Yet no controlled
experiment has measured whether offloading facts during pretraining
actually changes what the remaining capacity learns. We present the design
and pilot results of the first such experiment. Twin decoder models train
on byte-matched corpora under identical schedules; the only difference is
that the split arm sees every fact value wrapped in database lookup calls
whose value tokens are masked from the loss, so it is never rewarded for
memorizing a fact, only for asking a minimal external key-value store we
call the organizer. Because we generate the facts ourselves, the fact load
is a controlled, swept variable with known information content. Pilot runs
at 160M parameters establish the mechanism decisively and repeatedly: the
split model stores zero measured fact bits in its weights, answers 99.8
percent of fact probes with the organizer attached and none without it,
composes correct lookups for entities it has never seen in any training
document (at or near 100 percent across rounds), and roughly doubles its
dense twin on held-out fact-use question answering (56.9 to 70.6 percent
versus 27.1 to 32.9 percent) at identical parameter count, data content,
and data order. The pilots also produced a useful negative finding:
knowledge-free symbolic reasoning tasks did not rise above chance from
in-mixture exposure at pilot budgets, across three rounds of mixture and
difficulty remediation, which motivated a preregistered re-anchoring of
the primary endpoint to fact-use reasoning before any confirmatory run
began. The full preregistered battery, a dose-response sweep at 160M and a
billion-parameter paired confirmation, is running now on frozen margins
and decision rules. We argue the design closes the gap between a widely
deployed intuition and its first measurement, and that the pilot evidence
already justifies scaling the paradigm.
\end{abstract}

\section{Introduction}

A small language model has a fixed budget of parameters, and everything it
learns must share that budget. Facts about the world, which are numerous,
arbitrary, and individually low-value, compete for storage with the
reusable procedures that make a model useful: parsing, composing,
comparing, deducing. This tension is not hypothetical. The phi-3 technical
report describes deleting fact-heavy web pages from training data
precisely to ``leave more model capacity for reasoning for the mini size
models''~\cite{phi3}. A production model family was designed around the
belief that facts crowd out reasoning in small models.

That belief has never been tested directly. Data filtering, as in phi-3,
confounds fact removal with every other property of the removed documents.
Retrieval augmentation at inference time, the standard industrial remedy,
leaves the memorization pressure of pretraining fully intact, and a recent
systematic study finds that it leaves reasoning performance essentially
unchanged across model scales~\cite{memorize}. The Limited Memory Language
Model (LMLM) line of work~\cite{lmlm} built the missing training-time
mechanism, pretraining models that emit database lookup calls with the
retrieved values masked from the loss, and showed striking gains in
factual precision per parameter. On the question of whether the freed
capacity buys anything, however, the LMLM paper reports only parity on
five natural language understanding benchmarks and states in its
limitations that benchmarks disentangling knowledge from reasoning
``remain underexplored.'' The conjecture at the center of the small-model
design space is, in short, open.

This paper describes an experiment built to close it, and reports the
evidence from its pilot phase. The design principle throughout is that the
corpus itself is the measuring instrument. We generate the facts, so we
know exactly how many bits of factual content the dense arm is being asked
to store, we can sweep that load as a controlled dose, and we can annotate
the split arm's lookup calls with perfect coverage rather than with a
learned annotator. We generate the reasoning tasks, so held-out accuracy
is scored by construction-time oracles rather than heuristics, and
contamination is excluded by structure hashing rather than by n-gram
guesswork. Both arms of every pair see byte-identical natural text,
identical synthetic content, identical document order, and matched
per-component token budgets, so the contrast isolates a single bit of
information: whether the loss function rewards memorizing fact values.

Our contributions are threefold. First, a controlled experimental design,
to our knowledge the first, that combines training-time fact offloading, a
matched-pair baseline, a swept fact dose with known bit content, and
reasoning as a preregistered primary endpoint. Second, pilot-scale
evidence that the mechanism does exactly what the hypothesis requires,
replicated across three independent training rounds. Third, a
preregistered confirmatory battery, frozen before any confirmatory run
started, currently executing at 160M and one billion parameters.

\section{Related work}

\textbf{Training with external memory.} LMLM~\cite{lmlm} pretrains GPT-2
and Llama style models from scratch with entity facts wrapped in lookup
calls and masked from the loss, using a fine-tuned annotator model to
insert calls over Wikipedia text and a fuzzy retriever over 54.6M mined
triplets at inference. A 382M LMLM matches the factual precision of a 7B
baseline, and unplugging the database collapses it, establishing that
facts genuinely live outside the weights. We adopt this training mechanism
directly and gratefully. What LMLM does not do is measure a reasoning
consequence; its general-capability evaluation is five benchmarks near
their random floors at that scale, reported as parity. RETRO~\cite{retro}
trains with chunked cross-attention into a trillion-token store and gains
mainly on copying-adjacent perplexity and factual QA; its own analysis
attributes much of the gain to overlap exploitation. kNN-LM style post-hoc
memory actively degrades reasoning even with oracle
retrieval~\cite{knnlimits}. Memory layers~\cite{memorylayers} and related
sparse in-weight stores relocate memorization inside the gradient loop;
they test architectural separation, not external offloading, and their
stores cannot be edited or unplugged.

\textbf{Measurements adjacent to the hypothesis.} A 2026 scaling study of
retrieval-considerate pretraining~\cite{memorize} trades corpus size
against datastore size and finds retrieval substitutes for parameters on
knowledge tasks while reasoning tasks show minimal change; retrieval there
is inference-time, so weights still pay the memorization cost during
training. Charpentier et al.~\cite{moreroom} train small masked LMs with
idealized retrieval and observe reallocation toward syntax with a
regression in broader understanding, a caution that freed capacity does
not automatically become the capability one wants. Allen-Zhu and Li
quantify parametric fact storage at roughly two bits per parameter under
sufficient exposure~\cite{physics33} and supply the accounting method our
mediation analysis adapts, as well as the synthetic-task methodology our
reasoning corpora follow~\cite{igsm}.

\textbf{Theory.} Complementary learning systems theory
\cite{cls,clsupdate} holds that intelligent agents need a fast,
item-specific store alongside a slow structure learner, because arbitrary
item memory interferes with the extraction of generalizable structure.
Our experiment operationalizes the division of labor directly: the
organizer plays the fast store, the network is freed to be the slow
learner, and the theory's signature predictions (division of labor,
reduced interference, faster schema acquisition) map onto our H1 through
H4.

\section{Method}

\subsection{Two arms, one difference}

Every condition trains a pair of decoder-only transformers (RMSNorm,
rotary positions, SwiGLU, untied embeddings; 12 layers by 768 for the
160M class, 22 by 1792 for the 1B class) with the GPT-2 byte-pair
tokenizer extended by four special tokens. Members of a pair share
everything: initialization seed, data-order seed, batch schedule, learning
rate schedule, corpus content, document order, and per-component token
budgets matched within one percent. They differ in exactly one respect.
In the dense arm a biography sentence reads

\begin{quote}
\texttt{Kai Nakamura majored in Communications at ...}
\end{quote}

and every token carries loss. In the split arm the same sentence reads

\begin{quote}
\texttt{Kai Nakamura majored in <|db\_start|>Kai Nakamura,
major<|db\_retrieve|> Communications<|db\_end|> at ...}
\end{quote}

and the value tokens between the retrieve and end markers are excluded
from the loss. The split model conditions on values in context but
receives no gradient reward for predicting them. It is graded on knowing
when to open a call and how to compose the query. Nothing is retrieved
during training; the values sit inline, masked, so the training loop is
standard next-token prediction plus a label mask, with zero added FLOPs.

At inference the organizer comes alive. It is deliberately minimal: an
exact-match table from normalized \texttt{(entity, relation)} keys to
value strings, built by the corpus generator. A three-state interception
loop in the decoder watches for \texttt{<|db\_start|>}, records the query
the model composes, executes the lookup at \texttt{<|db\_retrieve|>}, and
force-decodes the returned value and closing marker through the model's
cache before generation resumes. Hits, misses, and malformed calls are
all counted. Using exact match over canonical synthetic keys removes
retriever quality as a confound entirely: if the split arm wins, no
reviewer can attribute the win to a clever retriever, because the store
is trivially perfect and its behavior fully audited.

\subsection{The corpus as instrument}

Both arms train on a five-component mixture (final frozen recipe: 54
percent natural web text from FineWeb-Edu, 23 percent synthetic
biographies, 12 percent synthetic mathematics, 8 percent synthetic
deduction, 3 percent fact-use question answering). The biographies are
the dose. Each of $N$ entities carries six attributes drawn from finite
pools with a combined information content of about 53 bits per entity,
rendered through more than twenty paraphrase templates per attribute,
eight sentence orderings, and three name forms. Sweeping $N$ (50k, 200k,
800k at 160M scale; 800k and 4M at 1B scale) sweeps the number of facts
the dense arm must store while the split arm may ignore. Because the
generator knows every fact, annotation coverage in the split arm is exact
by construction, unlike learned annotation pipelines which leak unwrapped
facts into ordinary text.

The reasoning components are knowledge-free by design. The mathematics
family follows the iGSM methodology~\cite{igsm}: random dependency graphs
over nonsense quantities (``crimson satchels in the Annex''), arithmetic
modulo 23 so no memorized number facts help, chain-of-thought solutions,
difficulty parameterized by operation count, and an independent oracle
that re-parses each prompt and recomputes the answer. The deduction family
poses Horn-clause chains over nonsense predicates with balanced yes and no
labels and a forward-chaining oracle. Fact-use QA asks questions whose
chain of thought must cite stored facts and then reason over them: date
comparisons, city equality checks, and extractions. Held-out sets are
disjoint by structure hash, and a separate probe set contains fresh
entities that exist only in the organizer, never in any training text,
which cleanly separates ``memorized the answer'' from ``learned the
lookup schema.''

\subsection{Endpoints, statistics, and gates}

The preregistration (frozen July 20, before any confirmatory run) fixes
the following. H1, primary: the split arm beats its dense twin on held-out
fact-use QA, split evaluated in its designed operating mode with the
organizer attached, dense closed-book; the confirmation contrast at 1B
uses a per-item paired bootstrap clustered by question template, with a
margin of $\max(2\hat\sigma_{\text{pool}}, 1.0)$ points and sign
consistency across seed pairs, where $\hat\sigma_{\text{pool}}$ is
estimated from the six 160M sweep pairs. H2, guardrail: split recall with
the organizer must be within two points of dense closed-book recall. H3,
mediation: split store-off recall below five percent and split fact bits
in weights below ten percent of the dense twin's, so any H1 delta
co-occurs with facts measurably living outside the weights. H4,
secondary: fewer tokens to fixed fact-use accuracy milestones. Synthetic
mathematics and deduction remain in every corpus as emergence-watch
secondaries with a pre-specified reporting trigger. A defensible null is
an explicit deliverable: a tight confidence interval around zero with all
guardrails intact.

Three gates protected compute before the battery. Gate C required the
lookup mechanics to work (they do, section 5.1). Gate B required the dose
to bind on the dense arm (it does, section 5.2). Gate A required the
knowledge-free tasks to be learnable at pilot budget; it failed three
times at about eight GPU-hours per round, a finding in its own right
(section 5.3), and triggered the endpoint re-anchoring above before any
confirmatory GPU-day was spent. We consider this sequence a feature of
the methodology worth emphasizing: every design risk was retired or
converted into a documented decision at pilot cost, and the confirmatory
battery launched with no unvalidated assumptions.

\section{Experimental setup}

Pilot runs train the 160M class for 0.8B tokens on a single NVIDIA L40S
each (about 130k tokens per second, roughly 1.7 hours per run). The
battery, running at the time of writing, comprises a dose-response sweep
(three loads, two arms, two seeds, 3.2B tokens per run) and a 1B paired
confirmation (two seed pairs, 10B tokens per run, dependency-chained
across scheduler wall-time limits), preceded by two short 1B calibration
runs whose selection rule is preregistered. All corpora were rebuilt on
the frozen recipe after the gate phase. The software stack is 124 unit
tests plus an end-to-end pipeline test; corpus builds are byte-identical
across serial and parallel execution; training is checkpoint-exact across
requeues; every run records its configuration, git state, and full
per-item evaluation outputs. Total cloud spend for everything reported
here: zero dollars, on a shared university cluster.

\section{Pilot results}

\subsection{The mechanism works, three times over}

Table~\ref{tab:mechanism} reports the split arm and its dense twin at the
200k-entity load for each of the three pilot rounds (which differ in
reasoning mixture share and difficulty band, giving an incidental
robustness check). The pattern is stable everywhere it should be and null
everywhere it should be.

\begin{table}[t]
\centering
\small
\begin{tabular}{lcccccc}
\toprule
& \multicolumn{2}{c}{Round 1} & \multicolumn{2}{c}{Round 2} & \multicolumn{2}{c}{Round 3}\\
\cmidrule(lr){2-3}\cmidrule(lr){4-5}\cmidrule(lr){6-7}
Measurement & dense & split & dense & split & dense & split\\
\midrule
Recall, store on (\%)        & --    & 99.8 & --   & 99.8 & --   & 99.4\\
Recall, store off (\%)       & 0.8$^{a}$  & 0.0  & 0.7$^{a}$ & 0.0  & 0.7$^{a}$ & 0.0\\
Fact bits in weights          & 40.3k & 0.0  & 25.3k & 0.0 & 21.3k & 0.0\\
Fact-use QA (\%)             & 32.9  & 70.6 & 27.1 & 56.9 & 31.2 & 58.1\\
Fresh-entity QA (\%)         & 0.2   & 100.0 & 0.4 & 100.0 & 0.4 & 99.6\\
Malformed lookups             & --    & 0    & --   & 0    & --   & 0\\
\bottomrule
\end{tabular}
\caption{Split arm versus dense twin, 200k entities, 0.8B-token pilots.
Rounds differ in reasoning share and difficulty (robustness check).
$^{a}$closed-book. Fresh-entity questions concern entities absent from
all training text; 500 items per round. Fact bits estimated from
attribute-level recall against pool entropies, following the accounting
of \cite{physics33}.}
\label{tab:mechanism}
\end{table}

Three results deserve emphasis. First, the separation is total: measured
fact storage in the split arm's weights is zero in every round, its
store-off recall is exactly zero, and during training its cross-entropy
on masked fact values stays near the uniform ceiling (about 9.5 nats
against a ceiling of 10.8) while its overall loss falls normally to about
2.4. The fast and slow stores are not a metaphor in this system; they are
a measured property of the trained model. Second, lookup is a learned,
generalizing skill: on entities the model has never read about, it
composes the correct canonical query at or near 100 percent, round after
round, with zero malformed calls in roughly ten thousand attempts.
Third, the behavioral payoff at the system level is large and stable: on
held-out fact-use QA the split arm scores between 27 and 38 points above
its dense twin, roughly doubling it, at identical size, content, and
order. This is precisely the operating mode the architecture is designed
for, and the margin persists across three independent trainings.

\subsection{The dose binds}

At fixed pilot budget the dense arm's closed-book recall falls
monotonically as the fact load rises: 1.1 percent at 50k entities, 0.8 at
200k, 0.3 at 800k. More facts strain the dense arm even at the pilot
scale, through the exposure-interference mechanism the design
anticipated, and the full-budget sweep quadruples exposures to size the
effect properly. The 1B calibration stage selects the dose that binds at
scale by a preregistered rule before the confirmation pair runs.

\subsection{A clean negative: in-mixture symbolic reasoning does not
ignite at pilot scale}

Gate A asked whether knowledge-free mathematics and deduction, present as
12 to 20 percent of the mixture, become executable at pilot budget. Three
rounds answered no (Table~\ref{tab:gatea}): accuracy stayed at the modular
arithmetic chance floor and the binary chance line respectively, through a
doubling of the reasoning share, a narrowed difficulty band, and finally a
floored curriculum in which the easiest problems require a single modular
operation. We read this as evidence that the arithmetic tables themselves
require drilling exposure well beyond a single pass at this scale, a
result consistent with the dedicated-corpus training regimes of the
synthetic reasoning literature and useful to anyone designing small-model
data mixtures. The gate system converted what could have been a wasted
550 GPU-hour battery into three eight-hour findings and a documented,
preregistered endpoint decision.

\begin{table}[t]
\centering
\small
\begin{tabular}{lcccc}
\toprule
Round (remediation) & iGSM dense & iGSM split & Deduction dense & Deduction split\\
\midrule
1 (share 12\%, op 2--8)            & 2.9 & 4.3 & 44.2 & 47.3\\
2 (share 20\%, op 2--6)            & 4.5 & 3.6 & 45.1 & 45.9\\
3 (floored: op 1--4, depth 1--2)   & 3.2 & 5.6 & 53.7 & 53.1\\
\midrule
Chance                              & \multicolumn{2}{c}{$\approx$4.3} & \multicolumn{2}{c}{50.0}\\
\bottomrule
\end{tabular}
\caption{Gate A across three pilot rounds (held-out accuracy, percent,
1{,}500 items per cell). Knowledge-free symbolic tasks stay at chance
under in-mixture exposure at 0.8B tokens and 160M parameters. Both tasks
remain in the battery corpora as emergence-watch secondaries at the 4x
full budget.}
\label{tab:gatea}
\end{table}

\subsection{Soundness summary}

We designed for auditability and believe the pilots demonstrate it.
Every contrast is a paired design isolating one causal bit. Every
synthetic answer is verified by an oracle that shares no code path with
generation. Held-out sets are disjoint by construction, not by filtering.
The store is exact-match, so retrieval quality cannot masquerade as
capacity. The mediation chain is measured at both ends: bits in weights
on one side, behavioral deltas on the other. The decision rules, margins,
variance estimators, exclusion policy, and endpoint hierarchy were frozen
and committed before any confirmatory run started, and the one endpoint
change this project has made was forced by pilot evidence, documented,
and completed before the battery launched. Negative and null outcomes
have preregistered, publishable forms.

\section{Discussion}

The pilot evidence establishes the premise of the memory-split hypothesis
at small scale: it is possible, with zero training-FLOP overhead and a
purely data-level intervention, to train a model whose factual knowledge
lives entirely in an external, auditable, instantly editable store, while
the model itself learns the general skill of consulting that store and
using what it returns. The behavioral dividend on reasoning-over-facts is
already visible and large in pilots. What remains, and what the running
battery will decide under frozen rules, is the quantitative question at
scale: how the dividend grows with fact load and whether it survives at a
billion parameters with tight intervals.

For the broader research program this experiment carries a second payoff
independent of the hypothesis test. The organizer, its training recipe,
its interception decoder, and the accounting instruments are exactly the
fast-memory substrate that a fast and slow architecture needs in
production: facts updatable without touching weights, deletion by row
removal, provenance by construction. The science and the platform
compound.

\textbf{Limitations.} Pilot results are at 160M parameters and 0.8B
tokens; the confirmatory numbers are pending and this paper will be
updated with them. The dose is synthetic; natural-corpus external
validity is a designed follow-up, not a claim we make here. The fact-use
endpoint compares the split system (model plus store) against a
closed-book dense model, which is the deployment-relevant contrast and
the preregistered one, but is a system-level rather than weights-only
comparison; the knowledge-free tasks that would isolate weights-only
capacity did not clear their learnability gate at pilot budget and are
tracked as secondaries at full budget. Incidental facts in the natural
bed are unmasked in both arms, which biases against the split arm and
makes the observed margins conservative. Seed replication at 1B is two
pairs, with variance estimated from six 160M pairs by a preregistered
pooling rule.

\section{Conclusion}

A hypothesis that already steers production small-model design has been
running ahead of its evidence. We built the experiment that measures it:
matched twins, a controlled fact dose, an exact and auditable external
memory, oracle-scored endpoints, and preregistered inference. The
mechanism is proven and replicated; the behavioral separation in the
architecture's operating mode is large and stable; the negative finding
en route is itself informative; and the confirmatory battery is running
on frozen rules with results due within days. We believe this paradigm,
validated at pilot scale for the cost of two engineer-days and zero cloud
dollars, is exactly the kind of instrument that merits scaling, and we
intend the next revision of this paper to carry the billion-parameter
answer.

\begin{thebibliography}{12}

\bibitem{phi3} M.~Abdin et al. Phi-3 technical report: A highly capable
language model locally on your phone. arXiv:2404.14219, 2024.

\bibitem{lmlm} L.~Zhao et al. Pre-training large memory language models
with internal and external knowledge. arXiv:2505.15962, 2025. ICLR 2026.

\bibitem{memorize} To memorize or to retrieve: Scaling laws for
RAG-considerate pretraining. arXiv:2604.00715, 2026.

\bibitem{moreroom} L.~Charpentier et al. More room for language:
Investigating the effect of retrieval on language models.
arXiv:2404.10939, 2024.

\bibitem{retro} S.~Borgeaud et al. Improving language models by
retrieving from trillions of tokens. arXiv:2112.04426, ICML 2022.

\bibitem{knnlimits} Great memory, shallow reasoning: Limits of kNN-LMs.
arXiv:2408.11815, NAACL 2025.

\bibitem{memorylayers} V.-P.~Berges, B.~O\u{g}uz, et al. Memory layers at
scale. arXiv:2412.09764, 2024.

\bibitem{physics33} Z.~Allen-Zhu and Y.~Li. Physics of language models:
Part 3.3, knowledge capacity scaling laws. arXiv:2404.05405, ICLR 2025.

\bibitem{igsm} Z.~Allen-Zhu and Y.~Li (and coauthors). Physics of
language models: Part 2.1, grade-school math and the hidden reasoning
process. arXiv:2407.20311, ICLR 2025.

\bibitem{cls} J.~L.~McClelland, B.~L.~McNaughton, and R.~C.~O'Reilly.
Why there are complementary learning systems in the hippocampus and
neocortex. Psychological Review, 102(3):419--457, 1995.

\bibitem{clsupdate} D.~Kumaran, D.~Hassabis, and J.~L.~McClelland. What
learning systems do intelligent agents need? Trends in Cognitive
Sciences, 20(7):512--534, 2016.

\end{thebibliography}

\end{document}
